\section*{Chapter \ref{ch:b2:precambrian-earth} --- The Precambrian Earth}

\begin{solution}{exo:b2:precambrian-earth:1}
(a) \omterm{def:b2:precambrian-earth:eons}{Hadean} $4567 - 4031 = \qty{536}{Myr}$; \omterm{def:b2:precambrian-earth:eons}{Archean} $4031 - 2500 = \qty{1531}{Myr}$;
\omterm{def:b2:precambrian-earth:eons}{Proterozoic} $2500 - 538.8 = \qty{1961}{Myr}$. (b) $(4567 - 538.8)/4567 = 0.882$: about seven
eighths. (c) $4567 - 4404 = \qty{163}{Myr}$ for the zircon; $4567 - 4030 = \qty{537}{Myr}$ for
the oldest rocks.
\end{solution}

\begin{solution}{exo:b2:precambrian-earth:2}
(a) The factors $e^{\lambda t}$ with $\lambda = \ln2/T$ and $t = \qty{3}{Gyr}$ are 1.593
(\ce{^{238}U}), 19.18 (\ce{^{235}U}), 1.160 (\ce{^{232}Th}) and 5.292 (\ce{^{40}K}): the
contributions were 3.053, 1.573, 2.424 and \qty{4.329}{pW}, \qty{11.38}{pW} in all, against
\qty{4.906}{pW} today. (b) $11.38/4.906 = 2.32$. (c) Thorium-232 leads today; potassium-40 led
then.
\end{solution}

\begin{solution}{exo:b2:precambrian-earth:3}
(a) $1616 + 18.0\times(28 - 32.2) = \qty{1540}{\celsius}$. (b) $32.2 + (1500 - 1616)/18.0 =
\qty{25.8}{\percent}$. (c) A whole-rock analysis is not the composition of a liquid: a
\omterm{def:b2:precambrian-earth:greenstone}{komatiite} flow can contain olivine crystals heaped up from the liquid, which add MgO. Its
\qty{38}{\percent} lies beyond the \qty{32}{\percent} of the richest \omterm{def:b2:precambrian-earth:greenstone}{komatiite} liquids, and
beyond the experiments on which the line rests.
\end{solution}

\begin{solution}{exo:b2:precambrian-earth:4}
Oldest rocks, \qty{4.03}{Ga}; first stromatolites, \qty{3.43}{Ga}; end of the Pilbara domes,
\qty{3.1}{Ga}; \omterm{def:b2:precambrian-earth:goe}{Great Oxidation Event}, about 2.4 to \qty{2.1}{Ga}; last common \omterm{def:b2:precambrian-earth:greenstone}{komatiites},
\qty{1.9}{Ga}; second rise of oxygen, about \qty{0.8}{Ga}; end of the Marinoan,
\qty{0.635}{Ga}; first Ediacaran fossils, about \qty{0.57}{Ga}.
\end{solution}

\begin{solution}{exo:b2:precambrian-earth:5}
(a) $\tfrac14\times1361\times0.706 = \qty{240.2}{W.m^{-2}}$ absorbed and $\sigma\times288^4 =
\qty{390.1}{W.m^{-2}}$, so $\varepsilon = 240.2/390.1 = 0.616$. (b) $\tfrac14\times1361\times0.4
= \qty{136.1}{W.m^{-2}}$, $T^4 = 136.1/(0.616\times\num{5.670e-8})$, $T = \qty{249.8}{K}$.
(c) At \qty{275.5}{K} the albedo is $0.6 - 0.306\times12.5/20 = 0.409$: $R = 340.25\times0.591 =
\qty{201.2}{W.m^{-2}}$ and $E = 0.616\,\sigma\times275.5^4 = \qty{201.2}{W.m^{-2}}$. (d) $R' =
340.25\times0.306/20 = \qty{5.21}{W.m^{-2}.K^{-1}}$, $E' = 4\varepsilon\sigma T^3 =
\qty{2.92}{W.m^{-2}.K^{-1}}$: $R' > E'$, so the equilibrium is unstable
(\cref{prop:b2:precambrian-earth:ebm}).
\end{solution}

\begin{solution}{exo:b2:precambrian-earth:6}
(a) $1/(1 + 0.4\times2.4/4.567) = 0.826$; $1/(1 + 0.4\times0.65/4.567) = 0.946$. (b) The warm
state needs at least 0.933 of today's sunlight: impossible \qty{2.4}{Ga} ago, possible
\qty{0.65}{Ga} ago. (c) $0.826/0.946 = 0.873$: \qty{13}{\percent} dimmer, close to the
\qty{10}{\percent} quoted, which comes from a different solar model.
\end{solution}

\begin{solution}{exo:b2:precambrian-earth:7}
(a) Sturtian $716.9 - 658.9 = \qty{58.0}{Myr}$; Marinoan $644.5 - 635.4 = \qty{9.1}{Myr}$;
interglacial $658.9 - 644.5 = \qty{14.4}{Myr}$. (b) $(58.0 + 9.1)/85 = 0.79$: about four fifths
of the period. (c) Its onset is only bracketed between 649.9 and \qty{639.0}{Ma}: on these
bounds the Marinoan lasted between about 3 and \qty{15}{Myr}.
\end{solution}

\begin{solution}{exo:b2:precambrian-earth:8}
(a) $\tau = \frac{0.05}{0.95}\times\frac{44.01}{28.97}\times\frac{\num{5.1e18}}{\num{0.26e12}} =
0.05263\times1.519\times\num{1.962e7} = \qty{1.57}{Myr}$. (b) $0.05/\num{278e-6} = 180$ times.
(c) Twice as long, \qty{3.14}{Myr}.
\end{solution}

\begin{solution}{exo:b2:precambrian-earth:9}
(a) The stones fell into soft mud and bent its laminae down: they are dropstones, released by
floating ice. (b) Ice reaching the sea \ang{10} from the equator is the expected trace of a
snowball, and very hard to obtain if only the polar regions are frozen. (c) A \omterm{def:b2:precambrian-earth:snowball}{cap carbonate}:
it records the sudden end of the glaciation, with intense weathering and the flooding of the
shelves by the meltwater.
\end{solution}

\begin{solution}{exo:b2:precambrian-earth:10}
(a) $0.94 > 0.933$, the warm tipping point: the warm equilibrium, $288\times0.94^{1/4} =
\qty{283.6}{K}$, lies just above \qty{283}{K}; and $0.94 < 1.227$, so the snowball also exists.
(b) $249.8\times0.94^{1/4} = \qty{246.0}{K}$. (c) $\varepsilon = 0.94\times1361\times0.4/(4\sigma
\times263^4) = 511.7/1085.1 = 0.472$. (d) $\sigma\times263^4 = \qty{271.3}{W.m^{-2}}$, and $(0.616
- 0.472)\times271.3 = \qty{39}{W.m^{-2}}$.
\end{solution}

\begin{solution}{exo:b2:precambrian-earth:11}
(a) $S_{\text{m}}/S_{\text{f}} = (T_1/T_2)^4\,(1 - a_{\text{w}})/(1 - a_{\text{i}})$: the
emissivity cancels. $(263/283)^4\times0.706/0.4 = 0.7459\times1.765 = 1.317$. (b) A smaller
contrast of albedo between ice and open water, such as darker, dusty ice, or a wider ramp of
temperature. (c) A snowball reflects so much sunlight that it stays cold under any sunlight or
greenhouse at which a warm Earth would freeze; to thaw it the forcing must exceed that of
freezing by the full width of the loop.
\end{solution}

\begin{solution}{exo:b2:precambrian-earth:12}
Near $T^*$, $C\dot\delta = -(E' - R')\delta$, so $\delta \propto e^{-t/\tau}$ with $\tau =
C/(E' - R')$. At the warm equilibrium $R' = 0$ and $E' = 4\times0.616\times\num{5.670e-8}\times
288^3 = \qty{3.34}{W.m^{-2}.K^{-1}}$; $C = 1025\times3992\times50 = \qty{2.05e8}{J.m^{-2}.K^{-1}}$,
so $\tau = \qty{6.1e7}{s}$, about \qty{1.9}{yr}. In this model the warm equilibrium keeps $R' =
0$ until it reaches \qty{283}{K}, so $\tau$ stays near two years up to the tipping point, then
the state vanishes. With a smoothly varying albedo, $R'$ grows as the equilibrium approaches the
fold, where $E' = R'$: $\tau$ grows without bound, a slowing down that announces the tipping
point.
\end{solution}

\begin{solution}{pb:b2:precambrian-earth:1}
\textbf{1.} $\num{0.0203e-6}\times0.9929\times\num{95.13e-6} = \qty{1.92}{pW.kg^{-1}}$
(\ce{^{238}U}); $\num{0.0203e-6}\times0.00711\times\num{568.47e-6} = \qty{0.0820}{pW.kg^{-1}}$
(\ce{^{235}U}); $\num{0.0795e-6}\times\num{26.28e-6} = \qty{2.09}{pW.kg^{-1}}$ (\ce{^{232}Th});
$\num{240e-6}\times\num{1.20e-4}\times\num{28.47e-6} = \qty{0.820}{pW.kg^{-1}}$ (\ce{^{40}K}).
\textbf{2.} \qty{4.91}{pW.kg^{-1}}.
\textbf{3.} $\ln2/T$: 0.155, 0.985, 0.0495 and \qty{0.555}{Gyr^{-1}}.
\textbf{4.} Factors 1.52, 14.3, 1.14 and 4.48: 2.91, 1.17, 2.39 and \qty{3.67}{pW.kg^{-1}}.
\textbf{5.} \qty{10.1}{pW.kg^{-1}}, 2.07 times today's.
\textbf{6.} Potassium-40. Its half-life, \qty{1.25}{Gyr}, is short: more than two half-lives
have passed since, and more than three quarters of the \ce{^{40}K} of that time has decayed.
\textbf{7.} $1/(1 + 0.4\times2.7/4.567) = 0.809$; $0.809\times1361 = \qty{1101}{W.m^{-2}}$.
\textbf{8.} $288\times0.809^{1/4} = \qty{273}{K}$, below \qty{283}{K}: not a valid warm state;
the Earth would freeze.
\textbf{9.} $(1101\times0.4/(4\times0.616\,\sigma))^{1/4} = \qty{237}{K}$.
\textbf{10.} $0.616\times0.809 = 0.498$.
\textbf{11.} $0.004/\num{278e-6} = 14$ to $0.03/\num{278e-6} = 108$ times.
\textbf{12.} The silicate-weathering thermostat: weathering slows in a cold climate, letting
volcanic carbon dioxide accumulate, and quickens in a warm one.
\textbf{13.} 0.826 of today's, below 0.933: no warm state; the Earth with today's air would be a
snowball.
\textbf{14.} $\frac{0.1}{0.9}\times1.519\times\num{1.962e7} = \qty{3.31}{Myr}$.
\textbf{15.} $\frac{0.2}{0.8}\times1.519\times\num{1.962e7} = \qty{7.45}{Myr}$.
\textbf{16.} The \omterm{def:b2:precambrian-earth:snowball}{cap carbonate}, a laminated dolostone sharply overlying the glacial deposits,
precipitated from a sea fed by intense weathering in the hot aftermath and flooded by
meltwater.
\textbf{17.} Iron(II) stays dissolved only in water without oxygen: the iron accumulates in the
anoxic ocean under the ice and precipitates as iron(III) where oxygen arrives, for instance with
meltwater.
\textbf{18.} $(1616 - 1450)/(32.2 - 23) = \qty{18.0}{\celsius}$ per per cent.
\textbf{19.} $1616 + 18.0\times(30 - 32.2) = \qty{1576}{\celsius}$.
\textbf{20.} \qty{1616}{\celsius}, the measured point.
\textbf{21.} $1616 - 18.0\times20.2 = \qty{1252}{\celsius}$.
\textbf{22.} $1616 - 20\times20.2 = \qty{1212}{\celsius}$: \qty{40}{\celsius} lower, a tenth of
the difference between the lavas; the conclusion does not depend on the choice.
\textbf{23.} \omterm{def:b2:precambrian-earth:greenstone}{Komatiites} came from plumes, hotter than the ambient mantle, and from large degrees of
melting deep down; a lava's temperature is not the potential temperature of its source. The
mantle as a whole was warmer by 200 to \qty{300}{K}.
\textbf{24.} The hottest \omterm{def:b2:precambrian-earth:greenstone}{komatiite} lava was about \textbf{\qty{364}{\celsius}} hotter than a
primary basalt of today's ridges.
\end{solution}
